\documentclass[11pt,a4paper]{article}
\usepackage[T1]{fontenc}
\usepackage[utf8]{inputenc}
\usepackage{lmodern}
\usepackage[margin=2.5cm,top=3cm]{geometry}
\usepackage{setspace}
\usepackage{parskip}
\setlength{\parskip}{0.7em}
\pagestyle{empty}

\begin{document}

\noindent
\textit{[To be printed on LSE letterhead and signed by Professor Albrecht Ritschl]}

\bigskip

\noindent [Date]

\bigskip

\noindent Dear Organising Committee,

I am writing to support the application of Aristeidis Grivokostopoulos for the rapporteur position at the Wealth of Nations at 250 conference at Glasgow.

Aristeidis joined the LSE Economic History PhD programme in September 2024, funded by an ESRC DTP studentship. He is a second-year PhD student with four working papers in circulation, conference presentations at WEHC (Lund, 2025) and the WEAI Annual Conference (San Francisco, 2025), and affiliations with the LSE Data Science Institute and the Hellenic Observatory.

His research is directly relevant to all three of his preferred assignments. His paper ``Between Civilisation and Constitution: Greek Parliamentary Debates on Gold Standard Adoption, 1875--1885'', prepared for the Mary Morgan Retirement Workshop (June 2026), applies Morgan's models as mediators framework to show how economic ideas fail when applied without their preconditions, connecting directly to the questions Besley and Ghatak raise about values and institutions and to Rodrik's argument about contextual versus universal policy. A companion paper uses 755 roll call votes to test whether fiscal shocks shift legislators from clientelist to programmatic behaviour, extending that institutional analysis to the legislative level. His work on debt restructuring in the Ottoman Empire, Egypt, and Greece uses bond-level London Stock Exchange data to estimate the credit market effects of conflict-driven fiscal supervision, with difference-in-differences estimates showing significant yield spread convergence in all three cases. This connects his empirical work directly to the fiscal infrastructure questions at the heart of Voth's chapter. He combines archival research with computational text analysis, writes well for audiences across economic history and the history of economic thought, and teaches EC311 (History of Economic Thought) at LSE.

I am happy to support his application.

\bigskip

\noindent Yours sincerely,

\bigskip\bigskip

\noindent Professor Albrecht Ritschl\\
Department of Economic History\\
London School of Economics and Political Science

\end{document}
